% GENERATED FILE. Edit canonical lessons, then run npm run generate:books.
% canonical-source-hash: fnv1a64:c5d8485498d32b06
% canonical-lessons: ES-C589-de-hecho, ES-C589-en-realidad, ES-C589-desde-luego, ES-C589-en-absoluto, ES-C589-en-el-fondo

\chapter{En realidad}
\label{ch:es-en-realidad}
\hlchaptermodality{589}

\begin{chapteropening}
\textbf{By the end of this chapter:} \emph{I can correct an idea, stress a point and say what is true deep down.}

Complete the last lesson of chapter 589: I can correct an idea, stress a point and say what is true deep down.
\end{chapteropening}

\section[de hecho]{de hecho — in fact}

\label{lesson:ES-C589-de-hecho}



\begin{quote}
Before the new one: say the Spanish for above all, then the Spanish for in general.
\end{quote}

\subsection*{You'll want to know: de hecho}
\textbf{de hecho} — \textquotedblleft{}in fact\textquotedblright{}.

\begin{quote}
De hecho, ya lo sabía.
\end{quote}

\begin{quote}
\emph{In fact, I already knew it.}
\end{quote}

It adds a stronger point or a correction.

\begin{cousinweb}
\emph{de hecho} is built on \emph{hecho}, inherited from Latin \emph{factus} 'done, made', the participle of \emph{facere} (EtymDB).
\end{cousinweb}

\begin{grammarlens}[title={What you've built}]
One more word for telling, explaining and describing at B1.
\end{grammarlens}

\subsection*{Your turn}
\begin{itemize}
\raggedright
  \item \emph{Say it:} \emph{de hecho}
  \item \emph{Say it:} \emph{de hecho}, once more
  \item \emph{Say it:} a sentence of your own with \emph{de hecho}
  \item \emph{Recall:} say the Spanish for above all, then the Spanish for in general, then say \emph{de hecho} again
\end{itemize}

\begin{tcolorbox}[breakable,colback=teal!4,colframe=teal!35!black,title={Before you move on}]
What does de hecho mean? (\textquotedblleft{}In fact\textquotedblright{}.) Say it once more.
\end{tcolorbox}
\section[en realidad]{en realidad — actually, really}

\label{lesson:ES-C589-en-realidad}



\begin{quote}
Before the new one: say the Spanish for in general, then the Spanish for in fact.
\end{quote}

\subsection*{You'll want to know: en realidad}
\textbf{en realidad} — \textquotedblleft{}actually, really\textquotedblright{}.

\begin{quote}
En realidad, no es tan caro.
\end{quote}

\begin{quote}
\emph{Actually, it is not so expensive.}
\end{quote}

It corrects what someone thinks.

\begin{cousinweb}
\emph{en realidad} is built on \emph{realidad}, a learned borrowing from Medieval Latin \emph{realitas}, built on \emph{res} 'thing' (EtymDB).
\end{cousinweb}

\begin{grammarlens}[title={What you've built}]
One more word for telling, explaining and describing at B1.
\end{grammarlens}

\subsection*{Your turn}
\begin{itemize}
\raggedright
  \item \emph{Say it:} \emph{en realidad}
  \item \emph{Say it:} \emph{en realidad}, once more
  \item \emph{Say it:} a sentence of your own with \emph{en realidad}
  \item \emph{Recall:} say the Spanish for in general, then the Spanish for in fact, then say \emph{en realidad} again
\end{itemize}

\begin{tcolorbox}[breakable,colback=teal!4,colframe=teal!35!black,title={Before you move on}]
What does en realidad mean? (\textquotedblleft{}Actually, really\textquotedblright{}.) Say it once more.
\end{tcolorbox}
\section[desde luego]{desde luego — certainly, of course}

\label{lesson:ES-C589-desde-luego}



\begin{quote}
Before the new one: say the Spanish for in fact, then the Spanish for actually.
\end{quote}

\subsection*{You'll want to know: desde luego}
\textbf{desde luego} — \textquotedblleft{}certainly, of course\textquotedblright{}.

\begin{quote}
Desde luego, te ayudo.
\end{quote}

\begin{quote}
\emph{Of course I will help you.}
\end{quote}

\emph{¡Desde luego!} can also show annoyance: honestly!

\begin{cousinweb}
\emph{desde luego} is built on \emph{luego}, inherited from Latin \emph{loco} 'on the spot', a form of \emph{locus} 'place' (EtymDB).
\end{cousinweb}

\begin{grammarlens}[title={What you've built}]
One more word for telling, explaining and describing at B1.
\end{grammarlens}

\subsection*{Your turn}
\begin{itemize}
\raggedright
  \item \emph{Say it:} \emph{desde luego}
  \item \emph{Say it:} \emph{desde luego}, once more
  \item \emph{Say it:} a sentence of your own with \emph{desde luego}
  \item \emph{Recall:} say the Spanish for in fact, then the Spanish for actually, then say \emph{desde luego} again
\end{itemize}

\begin{tcolorbox}[breakable,colback=teal!4,colframe=teal!35!black,title={Before you move on}]
What does desde luego mean? (\textquotedblleft{}Certainly, of course\textquotedblright{}.) Say it once more.
\end{tcolorbox}
\section[en absoluto]{en absoluto — not at all}

\label{lesson:ES-C589-en-absoluto}



\begin{quote}
Before the new one: say the Spanish for actually, then the Spanish for certainly.
\end{quote}

\subsection*{You'll want to know: en absoluto}
\textbf{en absoluto} — \textquotedblleft{}not at all\textquotedblright{}.

\begin{quote}
—¿Te molesta? —En absoluto.
\end{quote}

\begin{quote}
\emph{Does it bother you? Not at all.}
\end{quote}

Alone it means no, not at all, even without \emph{no}.

\begin{cousinweb}
\emph{en absoluto} is built on \emph{absoluto}, a learned borrowing from Latin \emph{absolutus} 'freed, complete', from \emph{absolvere} 'to set loose' (EtymDB).
\end{cousinweb}

\begin{grammarlens}[title={What you've built}]
One more word for telling, explaining and describing at B1.
\end{grammarlens}

\subsection*{Your turn}
\begin{itemize}
\raggedright
  \item \emph{Say it:} \emph{en absoluto}
  \item \emph{Say it:} \emph{en absoluto}, once more
  \item \emph{Say it:} a sentence of your own with \emph{en absoluto}
  \item \emph{Recall:} say the Spanish for actually, then the Spanish for certainly, then say \emph{en absoluto} again
\end{itemize}

\begin{tcolorbox}[breakable,colback=teal!4,colframe=teal!35!black,title={Before you move on}]
What does en absoluto mean? (\textquotedblleft{}Not at all\textquotedblright{}.) Say it once more.
\end{tcolorbox}
\section[en el fondo]{en el fondo — deep down, at heart}

\label{lesson:ES-C589-en-el-fondo}



\begin{quote}
Before the new one: say the Spanish for certainly, then the Spanish for not at all.
\end{quote}

\subsection*{You'll want to know: en el fondo}
\textbf{en el fondo} — \textquotedblleft{}deep down, at heart\textquotedblright{}.

\begin{quote}
En el fondo, es una buena persona.
\end{quote}

\begin{quote}
\emph{Deep down, he is a good person.}
\end{quote}

It gives what is true underneath the surface.

\begin{cousinweb}
\emph{en el fondo} is built on \emph{fondo} 'bottom', inherited from Latin \emph{fundus} 'bottom, base' (EtymDB).
\end{cousinweb}

\begin{grammarlens}[title={What you've built}]
One more word for telling, explaining and describing at B1.
\end{grammarlens}

\subsection*{Your turn}
\begin{itemize}
\raggedright
  \item \emph{Say it:} \emph{en el fondo}
  \item \emph{Say it:} \emph{en el fondo}, once more
  \item \emph{Say it:} a sentence of your own with \emph{en el fondo}
  \item \emph{Recall:} say the Spanish for certainly, then the Spanish for not at all, then say \emph{en el fondo} again
\end{itemize}

\begin{tcolorbox}[breakable,colback=teal!4,colframe=teal!35!black,title={Before you move on}]
What does en el fondo mean? (\textquotedblleft{}Deep down, at heart\textquotedblright{}.) Say it once more.
\end{tcolorbox}
